Because the device is intended for long-term unattended operation, the firmware must recover automatically from transient faults without any external intervention.
The dominant concern is not an incorrect measurement but a state in which the device remains awake at active current indefinitely, since such a state would drain the battery in a matter of days rather than years.
Two mechanisms address this, namely a per-cycle retry for failed transmissions and a no-progress watchdog that guarantees the device always returns to deep sleep.

The first mechanism handles the common case in which a single report is not delivered, for example because the network is not yet ready or because a transmission is interrupted or corrupted.
When the tracked transmission-complete callback reports a failure, the firmware does not enter sleep but instead schedules a short retry after a fixed interval of one second, after which the report is attempted again.
This turns a transient delivery failure into a brief, bounded retry loop rather than a lost measurement, and because each attempt is inexpensive the device recovers automatically once the link becomes available.

The second mechanism protects against the more serious case in which the active phase never completes at all, for example if a network join or the ZHA interview stalls and the device would otherwise stay awake waiting forever.
A no-progress watchdog records a deadline at start-up and the state machine checks on every iteration whether that deadline has passed.
If the active phase has not completed within the deadline, set to sixty seconds, the watchdog forces a short EM4 cycle that reboots the device and restarts the whole sequence from a clean state.
The deadline is armed once during initialization, and because every wake from EM4 is a full reset the watchdog is effectively re-armed at the start of every cycle without any explicit reset logic.
This guarantees that the device can never remain in the active phase for longer than the deadline, so a stalled join or interview costs at most one wasted cycle of active current rather than the entire battery.

This watchdog is implemented in software within the application rather than using the hardware watchdog peripheral of the EFR32MG22E, which is deliberately left disabled.
The hardware watchdog is designed to detect a hung main loop by requiring periodic servicing, whereas the failure mode of concern here is different, since the main loop continues to run normally while the network simply never reaches a usable state.
A software deadline measured against the application's own notion of progress matches this failure mode directly, and it also integrates cleanly with the deep-sleep entry path, since the watchdog requests the same EM4 cycle that the normal flow uses rather than triggering an independent hardware reset.

A further safeguard exists in the handling of the transmission-complete callback, which is delivered asynchronously by the stack.
A guard flag marks whether a tracked transmission is genuinely in flight, and any callback that arrives when no transmission is expected, such as a stale or duplicate completion, is rejected.
This prevents a late callback from being mistaken for the completion of a later cycle, which could otherwise allow the device to sleep before its current measurement had actually been delivered.